\documentclass[sigconf,nonacm]{acmart}

% --- clean up ACM boilerplate for a course report ---
\settopmatter{printacmref=false}   % no ACM reference-format block
\setcopyright{rightsretained}       % ignored under nonacm; silences the check
\renewcommand\footnotetextcopyrightpermission[1]{} % kill the permission footnote
\pagestyle{plain}

\usepackage{amsmath}
\usepackage{amssymb}

\begin{document}

\title{Certified but Compromised: Backdoor Attacks against Sampling-Based Proof-of-Inference}
\subtitle{Intermediate Report --- Trustworthy Machine Learning, Spring 2026}

\author{Ofek Lavi}
\affiliation{\institution{Tel Aviv University}\city{Tel Aviv}\country{Israel}}
\email{ofeklavi@mail.tau.ac.il}

\author{Erel Brazilay}
\affiliation{\institution{Tel Aviv University}\city{Tel Aviv}\country{Israel}}
\email{erelbrazilay@mail.tau.ac.il}

\author{Edo Kadosh}
\affiliation{\institution{Tel Aviv University}\city{Tel Aviv}\country{Israel}}
\email{edo@mail.tau.ac.il}

\begin{abstract}
Verifiable inference lets a client confirm that an answer from an untrusted cloud service
was produced by the intended model, without re-running the model locally. To avoid the
heavy prover cost of zkSNARK-based schemes, Anchuri et al.\ (SaTML 2026) recently proposed
a lightweight, sampling-based protocol: the prover commits to the inference \emph{execution
trace} and opens only a few entries along randomly sampled output-to-input paths, relying on
the assumption that \emph{functionally dissimilar} models yield \emph{representationally
dissimilar} traces. We observe that such schemes certify \emph{functional proximity}, not
\emph{model integrity}, and that backdoored models are the canonical counterexample:
identical to the intended model on the clean distribution, malicious only on rare
triggers---squarely in the scheme's blind spot. We will (i)~formalize an adversary
(committing a benign model while serving a backdoored one);
(ii)~empirically map evasion against the verifier's opening budget, trigger rarity, and
backdoor locality; and (iii)~attempt to prove that any inference-time verifier opening a sublinear
number of trace locations cannot detect sufficiently sparse backdoors except with negligible
probability. Our thesis: backdoor-robustness cannot be delivered at inference time by sublinear trace
sampling---it must come from verifiable training or model attestation.
\end{abstract}

\keywords{verifiable inference, proof-of-inference, backdoor attacks, trustworthy machine learning, property testing}

\maketitle

\section{Introduction}
Large models are increasingly served as opaque cloud APIs, and a client has no guarantee that
a returned answer came from the model the provider claims to run rather than a cheaper,
mis-trained, or tampered substitute; re-executing a billion-parameter model locally to check
is infeasible. \emph{Verifiable inference} has the prover accompany each answer with evidence
a verifier can check cheaply. Zero-knowledge approaches prove $y=M(x)$ for a bit-exactly
committed model~\cite{kang2022scaling,benno2026jolt} but incur prohibitive prover overhead. To
sidestep this, Anchuri et al.~\cite{anchuri2026verifiable} commit to the \emph{execution trace}
via vector commitments~\cite{catalano2013vector} and open only a few entries along
randomly sampled output-to-input paths, arguing soundness in a property-testing
style~\cite{goldreich2017property}: functionally dissimilar models yield representationally
dissimilar traces, detectable from very few opened locations.

We argue that this design admits a fundamental attack. Verifiable inference certifies that
\emph{a committed model was faithfully executed}; it says nothing about whether that model is
benign. Backdoors exploit exactly this gap: a stealthy backdoor is functionally \emph{identical}
to the intended model on the clean distribution and misbehaves only on rare
triggers~\cite{gu2019badnets,saha2020hidden}. For sampling-based schemes the gap is
\emph{structural}, not merely definitional: soundness is defined only up to functional
similarity, so a backdoor---a functional near-twin---satisfies the very predicate the scheme
certifies, and a sparse test that essentially never lands on a trigger cannot separate the two.

\noindent\textbf{Contributions.}
\begin{itemize}
  \item A precise account of the mismatch between what sampling-based proof-of-inference
  \emph{certifies} (functional proximity) and the integrity practitioners \emph{assume},
  instantiated as two adversaries against~\cite{anchuri2026verifiable}.
  \item An empirical demonstration of backdoored models \emph{accepted} by the verifier while
  achieving high trigger success, with a map of the evasion regime over the opening budget, trigger
  rarity, and backdoor locality.
  \item A lower bound showing that sublinear-budget trace sampling cannot detect sparse
  backdoors, implying backdoor-robustness belongs to verifiable \emph{training}, not inference.
\end{itemize}

\section{Related Work}
\textbf{Verifiable inference.} zkSNARK systems~\cite{kang2022scaling,benno2026jolt} prove
correctness against a bit-exactly committed model at high prover cost. Sampling-based and
refereed protocols trade exactness for efficiency; the scheme of Anchuri et
al.~\cite{anchuri2026verifiable} is the lightweight, statistically-grounded representative we
target. Our attack turns on a distinction: exact-binding SNARKs pin the model bit-for-bit, so
our evasion does not transfer to them.

\noindent\textbf{Backdoor and weight-tampering attacks.} Backdoors implant hidden behavior that
preserves clean-set accuracy and fires on a trigger~\cite{gu2019badnets}; hidden-trigger and
clean-label variants add stealth~\cite{saha2020hidden}.

\section{Work Plan}
We tackle the problem in five phases.

\noindent\textbf{Phase 0 --- Specification.} Extract from~\cite{anchuri2026verifiable} the exact
objects we must reproduce---trace generation, the vector commitment, the verifier's test, the
opening budget, and the path-sampling distribution---plus the formal (``other-model'') soundness
definition and its stated assumptions and limitations.

\noindent\textbf{Phase 1 --- Reproduction.} Reimplement prover and verifier on small, fully
instrumentable models (an MLP and a CNN on MNIST/CIFAR-10), and validate against the paper's expected behavior---accept an honest prover,
reject a functionally \emph{dissimilar} substitute---before attacking.

\noindent\textbf{Phase 2 --- Attacks.} \emph{(TM1, motivation)} the prover commits the backdoored
model and the verifier trusts that commitment as ``the intended model,'' yielding a
\emph{certified} malicious output and establishing that \emph{verified $\neq$ trustworthy}.
\emph{(TM2, main)} the prover commits a benign $M$ but serves a backdoored $M'$ that agrees with
$M$ on all non-trigger inputs; we measure verifier acceptance versus (a)~trigger frequency,
(b)~opened nodes/paths, and (c)~backdoor locality (a few-neuron trigger circuit vs.\ a distributed
one), and we attempt to \emph{localize} the backdoor in rarely-opened trace regions. Metrics for our success:
clean-accuracy gap $\Delta\mathrm{acc}$, trigger attack-success rate (ASR), verifier acceptance
under attack, and openings-to-detection curves.

\noindent\textbf{Phase 3 --- Lower bound.} We will attempt to prove that no scheme that knows a sublinear number of neuron weights can prove that the model is not backdoored w.h.p. 

\bibliographystyle{ACM-Reference-Format}
{\small\bibliography{references}}

\end{document}